\section{What the order slot reveals}
\label{sec:04_slot_beliefs}

The order slot gives the market one of three observations: a sale, zero, or a purchase.
This chapter derives the beliefs after these orders.
It then shows why the challenger can recover those beliefs from the price.
The proofs cover mixed profiles and wrong-signed orders. \paper{Sec.~3.2; Appendix~A.2}

\subsection{The order law under every profile}
\label{sec:04_slot_beliefs-law}

Write $\sigma_j(x)$ for the probability that the investor chooses $x$ in class $j$.
The three probabilities in each $\sigma_j$ sum to one.
At this stage, allow a sale in class $H$ and a purchase in class $L$.
The liquidity law is
\[
\eta(-1)=\beta,\qquad \eta(0)=1-2\beta,\qquad \eta(1)=\beta.
\]
An investor fills the slot with probability $\lambda$.
A liquidity trader fills it with probability $1-\lambda$.
These are alternative arrivals, so the conditional order probability is a sum:
\begin{align*}
f_j(x)&:=\Pr(X=x\mid j)\\
      &=(1-\lambda)\eta(x)+\lambda\sigma_j(x).
\end{align*}
Define the two liquidity masses
\[
b=(1-\lambda)\beta,\qquad b_0=(1-\lambda)(1-2\beta),\qquad u=\frac{\lambda}{b}.
\]
Thus $f_j(\pm1)=b+\lambda\sigma_j(\pm1)$ and $f_j(0)=b_0+\lambda\sigma_j(0)$.
The number $u$ compares informed presence with the liquidity mass in one direction.
It is not the posterior odds that an observed order is informed.

Because $0<\lambda<1$ and $0<\beta<1/2$, every liquidity mass is positive.
Hence $f_j(x)\ge (1-\lambda)\eta(x)>0$ for every order and class: every order occurs under every profile.
Bayes' rule therefore specifies every order posterior without an off-path belief.

The formulas do not ask whether the orders are optimal.
They apply to a profile with wrong-signed orders in both classes, and to a mixed profile.
\Cref{sec:benchmark-slot} works one such profile with exact fractions.

\subsection{Likelihood ratios and the posterior bounds}
\label{sec:04_slot_beliefs-bounds}

A bound on the likelihood ratio gives a bound on the posterior.

\begin{intuition}[Why one order cannot reveal the class]
A liquidity trader can place any order in either class.
The informed investor adds at most $\lambda$ to the probability of that order.
This limits the likelihood ratio, even under an extreme or mixed profile.
\end{intuition}

For a nonzero order, $0\le\sigma_j(x)\le1$ gives
\[
b\le f_j(x)\le b+\lambda.
\]
The largest ratio uses the largest numerator and the smallest denominator.
The smallest ratio reverses them:
\[
\frac{b}{b+\lambda}\le\frac{f_H(x)}{f_L(x)}\le\frac{b+\lambda}{b}.
\]
Since $\lambda/b=u$, these bounds are $1/(1+u)$ and $1+u$.
Bayes' rule gives $\mu(x)=\pi f_H(x)/[\pi f_H(x)+(1-\pi)f_L(x)]$ and $1-\mu(x)=(1-\pi)f_L(x)/[\pi f_H(x)+(1-\pi)f_L(x)]$.
Divide the first by the second; the denominators cancel:
\[
\frac{\mu(x)}{1-\mu(x)}=\frac{\pi}{1-\pi}\cdot\frac{f_H(x)}{f_L(x)}.\tag*{\paper{(A.3)}}
\]
This is the odds form of Bayes' rule.
The posterior rises strictly with these odds.
For a likelihood ratio $L$, its value is $\pi L/[\pi L+1-\pi]$.
Substitution of the two extreme ratios gives
\begin{align*}
m(\pi)&=\frac{\pi/(1+u)}{\pi/(1+u)+1-\pi}
       =\frac{\pi}{\pi+(1-\pi)(1+u)},\\
M(\pi)&=\frac{\pi(1+u)}{\pi(1+u)+1-\pi}
       =\frac{\pi(1+u)}{1+\pi u}.
\end{align*}
Thus $m(\pi)\le\mu(\pm1)\le M(\pi)$ under every profile. \paper{(4); Lemma~A.1(b)}
Both bounds lie strictly between zero and one.
Also $m(\pi)<\pi<M(\pi)$ because $u>0$.

\subsection{The zero order and the restriction \texorpdfstring{$\beta\le1/3$}{beta <= 1/3}}
\label{sec:04_slot_beliefs-zero}

At zero the relevant liquidity mass is $b_0$, rather than $b$.
The same argument gives
\[
b_0\le f_j(0)\le b_0+\lambda,\qquad
\frac{1}{1+\lambda/b_0}\le\frac{f_H(0)}{f_L(0)}\le1+\frac{\lambda}{b_0}.
\]
To use the previous bounds, we need
\begin{align*}
b_0\ge b
&\iff (1-\lambda)(1-2\beta)\ge(1-\lambda)\beta\\
&\iff 1-2\beta\ge\beta
\iff \beta\le\frac13.
\end{align*}
Then $\lambda/b_0\le\lambda/b=u$.
The zero-order likelihood ratio also lies between $1/(1+u)$ and $1+u$.
Hence $m(\pi)\le\mu(0)\le M(\pi)$. \paper{Lemma~A.1(c)}

\begin{caution}[A zero order can be the most informative order]
Do not treat zero as uninformative under an arbitrary profile.
It is uninformative under full orders because neither class chooses zero.
A one-sided profile can make zero good news.
\end{caution}

To see the failure when $\beta>1/3$, let class $H$ always abstain and class $L$ always sell.
Then $f_H(0)=b_0+\lambda$ and $f_L(0)=b_0$.
Since $b_0<b$,
\[
\frac{f_H(0)}{f_L(0)}=1+\frac{\lambda}{b_0}>1+\frac{\lambda}{b}=1+u.
\]
The zero-order likelihood ratio exceeds the top ratio, so $\mu(0)$ lies above $M(\pi)$.
\Cref{sec:benchmark-slot} gives an exact example.
This is a failure of the bound under a profile, regardless of investor optimality.
The nonzero-order bounds still hold for every $\beta<1/2$.

\subsection{Full orders and the slot beliefs}
\label{sec:04_slot_beliefs-full}

Under full orders, class $H$ buys and class $L$ sells with probability one.
The conditional laws and their ratios are
\[
\begin{array}{c|ccc}
x&-1&0&1\\ \hline
f_H(x)&b&b_0&b+\lambda\\
f_L(x)&b+\lambda&b_0&b\\
f_H(x)/f_L(x)&1/(1+u)&1&1+u\\
\mu(x)&m(\pi)&\pi&M(\pi)
\end{array}
\]
This establishes the three posteriors directly. \paper{Lemma~A.1(d)}
Full orders reach both bounds.
The zero order leaves the prior unchanged.

The slot is easier to describe by $u$ and $\beta$ than by $\lambda$ and $\beta$.
The presence probability then follows from them:
\begin{align*}
u&=\frac{\lambda}{(1-\lambda)\beta}
\iff u\beta(1-\lambda)=\lambda\\
&\iff \lambda=\frac{u\beta}{1+u\beta}.
\end{align*}
A larger $u$ moves both bounds away from the prior.
At the prior one half, $m=1/(2+u)$ and $M=(1+u)/(2+u)$, and the two bounds are symmetric about one half.
\Cref{sec:benchmark-slot} gives the benchmark slot and its beliefs.

\subsection{A price belief is an average, and the ceiling follows}
\label{sec:04_slot_beliefs-ceiling}

The challenger sees $P(X)$, rather than $X$.
For a price $\bar P$, let $A=\{x:P(x)=\bar P\}$.
Write $z_x=\pi f_H(x)+(1-\pi)f_L(x)$ for the unconditional probability of $x$.
Then
\begin{align*}
\bar\mu(\bar P)
&=\frac{\sum_{x\in A}\pi f_H(x)}{\sum_{x\in A}z_x}\\
&=\sum_{x\in A}\frac{z_x}{\sum_{y\in A}z_y}\,\mu(x).
\end{align*}
Every weight is positive, and the weights sum to one.
Therefore a price belief lies between the smallest and largest posteriors in its pool.
In particular, $\bar\mu(\bar P)\in[m(\pi),M(\pi)]$ when $\beta\le1/3$.
This argument needs no assumption that the price reveals the order.

\begin{resultbox}[The ceiling under every profile]
Let $\beta\le1/3$, $g_H\ge g_L$ and $B(M(\pi))<c_H$.
Then no price recruits a high-cost challenger under any profile, in or out of equilibrium, that has Bayes beliefs and optimal preparation.
\paper{Corollary~A.1}
\end{resultbox}

The proof has one inequality chain:
\[
B(\bar\mu(\bar P))=g_L+\bar\mu(\bar P)(g_H-g_L)
\le g_L+M(\pi)(g_H-g_L)<c_H.
\]
The first inequality uses $\bar\mu(\bar P)\le M(\pi)$ and $g_H-g_L\ge0$.
Thus high-cost preparation has a strictly negative net return at every price, also at a price that pools orders.
The comparison is strict, so the result holds under both challenger tie rules.
The universal ceiling uses the zero-order bound, so it needs $\beta\le1/3$.
Proposition~2 instead forces full orders first.
Its proof therefore needs only $\beta<1/2$.

\subsection{The entry floor and the competitive price}
\label{sec:04_slot_beliefs-price}

A zero-cost challenger prepares after every price.
Its gross profit is a nonnegative auction rent, so this choice is optimal.
The model fixes that choice even when its rent is zero.
If $\zeta(P)$ denotes high-cost preparation, entry after order $x$ is
\[
e(x)=\rho+(1-\rho)\zeta(P(x))\ge\rho>0.
\]
This is the entry floor. \paper{Lemma~A.2(a)}

Let $w_j=t_j-t_0$ and $w(\mu)=w_L+\DT\mu$.
The auction gives $w_H\ge w_L\ge0$.
Conditional on the class and the order, the cost draw alone determines preparation.
The cost is independent of both acquisition values and the class.
Also the order gives no further value information once the class is known.
Hence
\[
\E[V_T\mid j,X=x]=(1-e(x))t_0+e(x)t_j=t_0+e(x)w_j.
\]
Average this identity over the two classes:
\begin{align*}
P(x)=\E[V_T\mid X=x]
&=\mu(x)[t_0+e(x)w_H]+[1-\mu(x)][t_0+e(x)w_L]\\
&=t_0+e(x)\{w_L+\mu(x)(w_H-w_L)\}\\
&=t_0+e(x)[t_L-t_0+\DT\mu(x)].
\end{align*}
This is the competitive price. \paper{(5); (A.4); Lemma~A.2(b)}
It includes the expected proceeds from the entry that the price itself causes.

\subsection{Price sufficiency, the threshold, and strict price order}
\label{sec:04_slot_beliefs-sufficiency}

Assume $\DT>0$ and $g_H>g_L$.

\begin{intuition}[Why the circularity is harmless]
The price depends on entry, and entry depends on the price.
You do not need to solve this circle to invert the price.
Two orders with one price have one entry probability, and that entry is at least $\rho$.
The price formula then leaves only the posterior to differ, and it cannot.
\end{intuition}

First prove price sufficiency without an assumption about entry as a function of the order posterior.
If $P(x)=P(x')=\bar P$, the challenger faces the same price and cost law.
Thus $e(x)=e(x')=:\bar e$.
The price formula gives
\[
0=P(x)-P(x')=\bar e\DT[\mu(x)-\mu(x')].
\]
Since $\bar e\DT\ge\rho\DT>0$, the two posteriors are equal.
The price belief averages equal posteriors, so
\[
\bar\mu(P(x))=\mu(x).
\]
Equivalently, the challenger can recover the posterior from
\[
\mu(x)=\frac{P(x)-t_0-e(P(x))(t_L-t_0)}{e(P(x))\DT}.
\]
Here $e(P)$ denotes the entry probability associated with price $P$. \paper{Lemma~A.2(c)}
The denominator is positive because of the entry floor.

Since $g_H>g_L$, $B(\mu)\ge c_H$ holds exactly when $\mu\ge\tau=(c_H-g_L)/(g_H-g_L)$, the threshold of \Cref{sec:03_auction-derivatives}.
By sufficiency, the challenger's belief at $P(x)$ is $\mu(x)$.
Under rule C1, an indifferent high-cost challenger prepares.
Thus entry is $\rho$ below $\tau$ and one at or above $\tau$.
Under rule CM, entry at $\tau$ can be any value in $[\rho,1]$.
The threshold can lie outside $[0,1]$; the same inequalities still apply.

Take $\mu_2>\mu_1$, and write $e_i$ for the associated entry.
Optimal preparation gives $e_2\ge e_1$ under either rule.
Add and subtract $e_2w(\mu_1)$:
\begin{align*}
P_2-P_1
&=e_2w(\mu_2)-e_1w(\mu_1)\\
&=e_2[w(\mu_2)-w(\mu_1)]+(e_2-e_1)w(\mu_1)\\
&=e_2\DT(\mu_2-\mu_1)+(e_2-e_1)[w_L+\DT\mu_1]\\
&\ge\rho\DT(\mu_2-\mu_1)>0.
\end{align*}
The second term is nonnegative because $e_2\ge e_1$ and $w_L+\DT\mu_1\ge0$.
Thus prices rise strictly with the posterior. \paper{(A.5); Lemma~A.2(e)}
Equal prices can pool equal posteriors, as in the weak economy.
They cannot pool distinct posteriors under these assumptions.

Under full orders in the strong economy, only the purchase posterior $M(\pi)$ exceeds $\tau$.
A purchase therefore recruits the high-cost challenger, and a sale or zero does not.
Entry is one after a purchase and $\rho$ after the other two orders.
The three prices are distinct, and the inverse formula recovers each posterior.
\Cref{sec:benchmark-slot} gives the three prices at the benchmark.

\begin{panelnote}[What changed from the Laplace version]
The Laplace version bounded the likelihood ratio through the noise density; the slot bounds it through the liquidity mass in each order.
The benchmark slot is chosen so that its top posterior equals the Laplace value.
\Cref{sec:10_robustness-sizes} derives this match and explains why it concerns the bound only.
The zero order is new: it needs $\beta\le1/3$.
\end{panelnote}
